\chapter{Revisão bibliográfica}
\label{chap:revisao}

\section{Projeto de propellers}
\label{sec:projeto_propellers}
\todo{Contextualizar o processo de projeto e as principais grandezas de desempenho. O livro de Gur pode servir como referência organizadora desta seção.}

\section{Aerodinâmica de propellers}
\label{sec:aerodinamica}

\subsection{Teoria da quantidade de movimento}
\todo{Apresentar apenas a teoria necessária para construir a BEMT.}

\subsection{Teoria do elemento de pá}
\todo{Definir velocidades, ângulos, forças locais e integração radial.}

\subsection{Blade Element Momentum Theory (BEMT)}
\label{sec:bemt}
\todo{Derivar a formulação adotada, hipóteses, procedimento iterativo e correções utilizadas no código.}

\subsection{Polares aerodinâmicas e XFOIL}
\label{sec:polares}
\todo{Documentar origem/interpolação de $C_L$ e $C_D$, dependência com Reynolds e eventuais limites do XFOIL.}

\section{Modelagem aeroacústica de propellers}
\label{sec:aeroacustica}

\subsection{Ruído tonal e frequência de passagem de pá}
\todo{Apresentar os mecanismos relevantes ao escopo e definir BPF e harmônicos.}

\subsection{Formulação de Hanson}
\label{sec:hanson}
\todo{Apresentar a formulação efetivamente utilizada no código, suas entradas, hipóteses e saídas.}

\section{Modelagem estrutural de pás}
\label{sec:estrutural_teoria}
\todo{Definir o modelo analítico estrutural, esforços considerados, tensões, deslocamentos, torção e critérios de integridade.}

\section{Otimização multidisciplinar de propellers}
\label{sec:otimizacao_teoria}
\todo{Revisar variáveis de projeto, funções objetivo, restrições e métodos de otimização adequados ao problema. A escolha final do algoritmo deve ser justificada mais adiante.}

\section{Síntese da revisão bibliográfica}
\label{sec:sintese_revisao}
\todo{Comparar os trabalhos mais relevantes e explicitar a lacuna que o presente TCC pretende atender. Uma tabela comparativa pode ser útil.}
